\section*{Bab \ref{ch:g11:plate-motions} --- Lempeng dan Geraknya}

\begin{solution}{exo:g11:plate-motions:1}
\omterm{def:g11:plate-motions:rigid-plate}{Lempeng kaku} adalah potongan \omterm{def:g11:lithosphere-asthenosphere:lithosphere}{litosfer} yang bergerak sebagai satu blok: jarak
antara titik-titiknya tidak berubah, dan deformasinya terkumpul di zona-zona
sempit di sepanjang batasnya. Gempa bumi menunjukkannya: hampir semuanya terjadi
di sepanjang sabuk-sabuk yang sempit, yaitu batas lempeng, dan sangat sedikit
yang terjadi di dalam lempeng.
\end{solution}

\begin{solution}{exo:g11:plate-motions:2}
\qty{72}{mm} setahun sama dengan \qty{72}{km} per sejuta tahun. Dalam 30 tahun
stasiun itu bergerak $72 \times 30 = \qty{2160}{mm}$, sekitar \qty{2.2}{m};
dalam sejuta tahun, \qty{72}{km}.
\end{solution}

\begin{solution}{exo:g11:plate-motions:3}
$v_E = 30/10 = \qty{3}{cm}$ setahun dan $v_N = 40/10 = \qty{4}{cm}$ setahun.
Kecepatannya $\sqrt{3^2 + 4^2} = \qty{5}{cm}$ setahun. Sudut $\beta$ dengan
arah utara memenuhi $\tan\beta = 3/4$, jadi $\beta \approx \ang{37}$: stasiun
itu bergerak ke arah timur laut, \ang{37} di timur arah utara.
\end{solution}

\begin{solution}{exo:g11:plate-motions:4}
\omterm{def:g11:plate-motions:hotspot-track}{Jejak titik panas} adalah deretan gunung api dan gunung bawah laut yang dibawa
lempeng menjauhi titik panas ketika lempeng itu bergerak di atasnya, makin tua
seiring jaraknya dari gunung api yang aktif. Gunung api yang paling muda,
Kilauea, berdiri di ujung tenggara rantai itu, di Pulau Hawaii; umurnya
bertambah ke arah barat-barat laut. Lempeng membawa gunung-gunung apinya
menjauhi titik panas searah dengan geraknya: lempeng Pasifik bergerak ke arah
barat-barat laut.
\end{solution}

\begin{solution}{exo:g11:plate-motions:5}
$519/5.1 \approx \qty{102}{km}$ per sejuta tahun, yaitu \qty{102}{mm}
setahun, sekitar \qty{10}{cm} setahun.
\end{solution}

\begin{solution}{exo:g11:plate-motions:6}
Komponennya $v_E = 63/31 \approx \qty{20.3}{mm}$ setahun dan $v_N = 48/31 \approx
\qty{15.5}{mm}$ setahun. Kecepatannya $\sqrt{20.3^2 + 15.5^2} \approx
\qty{25.5}{mm}$ setahun; $\tan\alpha = 15.5/20.3 \approx 0.76$ memberikan $\alpha
\approx \ang{37}$: stasiun itu bergerak \ang{37} di utara arah timur, ke arah
timur laut. Mauna Kea bergerak $72/25.5 \approx 2.8$ kali lebih cepat.
\end{solution}

\begin{solution}{exo:g11:plate-motions:7}
(1) Dalam milimeter setahun,
\[
  \vec v_{\mathrm{NA/EU}} = (-10.5 - 7.7,\ 20.6 - 16.3) = (-18.2, +4.3).
\] Kecepatannya $\sqrt{18.2^2 + 4.3^2} \approx
\qty{18.7}{mm}$ setahun, dan $\tan\alpha = 4.3/18.2 \approx 0.24$ memberikan
\ang{13} di utara arah barat. (2) Kedua lempeng saling menjauh melintasi
punggung samudra dengan kecepatan sekitar \qty{1.9}{cm} setahun: inilah \omterm{def:g11:drift-to-spreading:spreading-rate}{laju
pemekaran penuh} punggung itu di Islandia. (3) Nilai itu sesuai dengan
\qty{2}{cm} setahun yang disebut untuk Atlantik Utara.
\end{solution}

\begin{solution}{exo:g11:plate-motions:8}
(1) $v = \pi \times 0.651 \times 6371/180 \approx \qty{72.4}{km}$ per sejuta
tahun, sekitar \qty{7.2}{cm} setahun. (2) $72.4 \times \sin\ang{30} \approx
\qty{36.2}{km}$ per sejuta tahun, sekitar \qty{3.6}{cm} setahun. (3) Di kutub
rotasi, jika kutub itu terletak pada lempeng. Kecepatannya berubah dari titik ke
titik, dari nol di kutub sampai nilai terbesarnya pada \ang{90} dari kutub:
seluruh lempeng tidak mempunyai satu kecepatan saja, hanya kecepatan di setiap
titik (atau rata-rata atas permukaannya).
\end{solution}

\begin{solution}{exo:g11:plate-motions:9}
(1) Dalam milimeter setahun,
\[
  \vec v_{\mathrm{IN/EU}} = (40.6 - 24.0,\ 38.6 + 3.3) = (+16.6, +41.9).
\] Kecepatannya $\sqrt{16.6^2 + 41.9^2} \approx
\qty{45}{mm}$ setahun; sudut $\beta$ dengan arah utara memenuhi $\tan\beta =
16.6/41.9 \approx 0.40$, jadi $\beta \approx \ang{22}$: India bergerak ke arah
Eurasia, \ang{22} di timur arah utara. (2) \qty{45}{km} per sejuta tahun, jadi
\qty{450}{km} dalam 10 juta tahun.
\end{solution}

\begin{solution}{exo:g11:plate-motions:10}
$(2432 - 1058)/(27.7 - 10.3) = 1374/17.4 \approx \qty{79}{km}$ per sejuta
tahun, sekitar \qty{7.9}{cm} setahun, sedikit kurang dari \qty{9.1}{cm} setahun.
Umur masing-masing gunung api mempunyai ketidakpastian sekitar setengah juta
tahun; selisih umurnya, $17.4 \pm 1.0$ juta tahun, memberikan kecepatan antara
sekitar 75 dan \qty{84}{km} per sejuta tahun. Di antara kedua gunung api ini
lempeng tampaknya bergerak sedikit lebih lambat daripada rata-ratanya.
\end{solution}

\begin{solution}{exo:g11:plate-motions:11}
(1) $\vec v_{A/C} = \vec v_{A/B} + \vec v_{B/C}$: \qty{3}{cm} setahun ke timur
ditambah \qty{4}{cm} setahun ke utara. Kedua panah itu membentuk sudut siku-siku,
dan sisi ketiga segitiga itu panjangnya $\sqrt{3^2 + 4^2} = \qty{5}{cm}$ setahun,
ke arah \ang{37} di timur arah utara ($\tan\beta = 3/4$). (2) $\vec v_{C/A}
= -\vec v_{A/C}$: \qty{5}{cm} setahun ke arah yang berlawanan, \ang{37} di barat
arah selatan.
\end{solution}

\begin{solution}{exo:g11:plate-motions:12}
(1) $54/0.375 = \qty{144}{km}$ per sejuta tahun, \qty{14.4}{cm} setahun. (2)
Dari $29/0.375 \approx 77$ sampai $79/0.375 \approx \qty{211}{km}$ per sejuta
tahun: antara sekitar 8 dan \qty{21}{cm} setahun. (3) Di dekat titik panas,
jaraknya tidak lebih besar daripada ukuran sebuah gunung api atau wilayah tempat
titik panas memasok gunung api, dan umurnya tidak lebih besar daripada waktu yang
diperlukan sebuah gunung api untuk tumbuh: keduanya sangat tidak pasti secara
relatif. Jauh dari titik panas, kesalahan \qty{25}{km} pada \qty{2400}{km}, atau
setengah juta tahun pada 28 juta tahun, hanya satu atau dua persen.
\end{solution}

\begin{solution}{exo:g11:plate-motions:13}
(1) Sepuluh kali sebaran itu sama dengan \qty{50}{mm}: $50/72 \approx 0.7$
tahun, sekitar delapan bulan. (2) $50/5 = 10$ tahun. (3) Dalam jutaan tahun
lempeng-lempeng bergerak puluhan sampai ratusan kilometer, jarak yang direkam
oleh jalur-jalur magnetik dan yang mudah diukur di laut; penentuan posisi
dengan ketelitian beberapa milimeter, yang membuat gerak tampak dalam beberapa
bulan, baru muncul jauh kemudian.
\end{solution}

\begin{solution}{exo:g11:plate-motions:14}
(1) Dalam milimeter setahun,
\[
  \vec v_{\mathrm{SA/AF}} = (-4.9 - 27.1,\ 11.0 - 17.9) = (-32.0, -6.9).
\] Kecepatannya $\sqrt{32.0^2 + 6.9^2} \approx
\qty{32.7}{mm}$ setahun, dan $\tan\alpha = 6.9/32.0 \approx 0.22$ memberikan
\ang{12} di selatan arah barat. (2) \omterm{def:g11:drift-to-spreading:spreading-rate}{Laju pemekaran penuh} sekitar \qty{3.3}{cm}
setahun, laju setengahnya sekitar \qty{1.6}{cm} setahun. (3) Laju setengah
selama beberapa juta tahun terakhir, \qty{1.6}{cm} setahun, sedikit lebih kecil
daripada \qty{1.9}{cm} setahun yang dirata-ratakan selama 66 juta tahun:
Atlantik Selatan mungkin membuka sedikit lebih lambat akhir-akhir ini daripada
rata-ratanya, atau perbedaan itu mungkin berasal dari ketidakpastian umur
sedimen. Kedua pengukuran sepakat bahwa punggung itu lambat.
\end{solution}

\begin{solution}{exo:g11:plate-motions:15}
(1) Garis lurus yang naik \qty{20}{mm} setahun, penurunan mendadak sebesar
\qty{50}{mm} setelah lima tahun, lalu garis yang naik lagi dengan kemiringan
yang sama. (2) Stasiun itu berakhir $10 \times 20 - 50 = \qty{150}{mm}$ di timur
titik awalnya: $150/10 = \qty{15}{mm}$ setahun, seperempat terlalu kecil.
(3) Setiap paruh memberikan \qty{20}{mm} setahun. (4) Lompatan itu adalah gerak
mendadak tersendiri yang disebabkan oleh gempa bumi: ahli \omterm{def:g10:measuring-the-earth:geodesy}{geodesi} mengukur
besarnya dari posisi tepat sebelum dan tepat sesudahnya, menghilangkannya dari
deret itu (atau mencocokkan kedua bagian dengan kemiringan yang sama dan satu
undakan di antaranya), dan mempertahankan \qty{20}{mm} setahun sebagai kecepatan
stasiun itu.
\end{solution}

\begin{solution}{pb:g11:plate-motions:1}
\textbf{1.} Mauna Kea (\num{0.375} juta tahun) yang paling muda, Suiko
(tengah) (\num{64.7} juta tahun) yang paling tua. Umurnya bertambah seiring
jarak dari Kilauea.

\textbf{2.} $\qty{1}{km} = 10^6~\unit{mm}$ dan sejuta tahun sama dengan $10^6$
tahun: \qty{1}{km} per sejuta tahun sama dengan $10^6~\unit{mm}/10^6$ tahun,
\qty{1}{mm} setahun.

\textbf{3.} $519/5.1 \approx \qty{102}{mm}$ setahun.

\textbf{4.} Necker $1058/10.3 \approx 102.7$; Laysan $1818/19.9 \approx
91.4$; Midway $2432/27.7 \approx 87.8$; Colahan $3128/38.6 \approx 81.0$;
Abbott $3280/38.7 \approx 84.8$ (milimeter setahun).

\textbf{5.} Titik-titik (54; 0.375), (519; 5.1), (780; 7.2), (1058; 10.3),
(1818; 19.9), (2432; 27.7), (3128; 38.6), (3280; 38.7) terletak dekat dengan
garis lurus melalui titik asal, dengan sedikit sebaran (terhadap garis
\qty{91}{km} per sejuta tahun, Kauai, Nihoa, dan Necker terletak sedikit di
bawahnya, Laysan tepat padanya, Colahan dan Abbott sedikit di atasnya).

\textbf{6.} Setiap gunung api dibangun di atas titik panas dan sejak itu
bergerak bersama lempeng. Dalam waktu $t$, lempeng telah bergerak sejauh $v\,t$
ke arah yang tetap, sehingga gunung api itu terletak pada jarak $d = v\,t$ dari
titik yang sekarang berada di atas titik panas: $t = d/v$.

\textbf{7.} $(102.7 + 91.4 + 87.8 + 81.0 + 84.8)/5 \approx \qty{89.5}{mm}$
setahun.

\textbf{8.} 89.5 dibandingkan dengan \qty{91}{km} per sejuta tahun: selisihnya
kurang dari \qty{2}{\percent}.

\textbf{9.} $519/91 \approx 5.7$ juta tahun, dibandingkan dengan umur \num{5.1}
juta tahun: dekat, dengan Kauai sedikit lebih muda daripada yang diramalkan garis
rata-rata itu.

\textbf{10.} Gunung-gunung api makin tua ke arah barat-barat laut: lempeng
membawanya ke arah itu. Lempeng Pasifik bergerak ke arah barat-barat laut di
atas titik panas.

\textbf{11.} $(4860 - 3493)/(64.7 - 42.4) = 1367/22.3 \approx \qty{61}{km}$
per sejuta tahun, sekitar \qty{6.1}{cm} setahun.

\textbf{12.} $65/91 \approx 0.71$: di sepanjang rantai Emperor lempeng bergerak
sekitar \qty{30}{\percent} lebih lambat daripada di sepanjang rantai Hawaii.

\textbf{13.} $3668/50 \approx \qty{73}{km}$ per sejuta tahun, sekitar
\qty{7.3}{cm} setahun, kurang dari \qty{91}{km} per sejuta tahun. Garis melalui
titik asal memakai umur tahun 1987, yang menempatkan ujung rantai Hawaii pada 39
sampai 43 juta tahun; umur yang lebih baru, sekitar 50 juta tahun, untuk awal
tikungan menyebarkan jarak yang sama atas waktu yang lebih panjang. Keduanya
benar jika lempeng bergerak lebih lambat pada bagian awal rantai Hawaii dan lebih
cepat kemudian: kecepatannya tidak tetap.

\textbf{14.} Perubahan arah gerak lempeng Pasifik, sekitar 50 juta tahun yang
lalu.

\textbf{15.} Arah rantai adalah arah gerak lempeng; jarak antarumur di
sepanjang rantai itu adalah kecepatannya. Sebuah lempeng dapat bertambah cepat
atau melambat tanpa berbelok, atau berbelok tanpa mengubah kecepatan, sehingga
keduanya saling bebas. Perubahan kecepatan tanpa perubahan arah di sepanjang
rantai Emperor menunjukkan bahwa rekaman kecepatannya dan tikungan itu adalah
peristiwa yang terpisah: tikungan itu tidak dibuat oleh titik panas yang
melambat.

\textbf{16.} $\sqrt{63.2^2 + 34.5^2} \approx \qty{72.0}{mm}$ setahun, dan
$\tan\alpha = 34.5/63.2 \approx 0.55$ memberikan \ang{29} di utara arah barat.

\textbf{17.} Barat-barat laut adalah \ang{22.5} di utara arah barat, barat laut
\ang{45}: arah stasiun itu, \ang{29} di utara arah barat, terletak di antara
keduanya, dekat dengan arah rantai itu. Arahnya sesuai.

\textbf{18.} Kecepatan stasiun diukur terhadap kerangka acuan yang di dalamnya
\omterm{def:g11:lithosphere-asthenosphere:lithosphere}{litosfer} secara keseluruhan tidak berputar, sedangkan rantai itu memberikan gerak
terhadap titik panas: jika kedua kerangka acuan itu bergerak relatif satu sama
lain, kecepatannya berbeda. Dan stasiun itu mengukur 30 tahun terakhir, rantai
itu 40 juta tahun terakhir, yang selama itu kecepatannya mungkin telah berubah.
Sebaran umur juga meninggalkan ketidakpastian beberapa milimeter setahun pada
kecepatan rantai itu.

\textbf{19.} $100/91 \approx 1.1$ juta tahun.

\textbf{20.} \textbf{Lempeng Pasifik telah bergerak di atas titik panas Hawaii
dengan kecepatan sekitar \qty{9}{cm} setahun} selama 40 juta tahun terakhir, ke
arah barat-barat laut; sebelum tikungan, lempeng itu bergerak lebih lambat, dan
ke arah yang lain.
\end{solution}
